% ============================================================================
%  El problema gravitatorio de N cuerpos — informe (español)
%  Compilación: python tools/build_documents.py   (necesita xelatex)
% ============================================================================
\documentclass[11pt,a4paper]{article}

\usepackage{amsmath,amssymb}
\usepackage{fontspec}
\usepackage{unicode-math}
\setmathfont{Latin Modern Math}
\usepackage{polyglossia}
\setdefaultlanguage{spanish}
\usepackage{etoolbox}
\gappto\captionsspanish{\renewcommand{\tablename}{Tabla}}
\usepackage[margin=2.3cm]{geometry}
\usepackage{graphicx}
\usepackage{booktabs}
\usepackage{siunitx}
\sisetup{exponent-product=\times, output-decimal-marker={,}, group-minimum-digits=5, group-digits=integer}
\usepackage[font=small,labelfont=bf,margin=0.5cm]{caption}
\usepackage{float}
\usepackage{xcolor}
\definecolor{link}{HTML}{1f5fae}
\usepackage[colorlinks=true,linkcolor=link,citecolor=link,urlcolor=link]{hyperref}
\hypersetup{pdftitle={Jesús García Izquierdo - https://github.com/jesus-garcia-izquierdo/N-body-problem}, pdfauthor={Jesús García Izquierdo}}
\usepackage{enumitem}
\setlist{nosep,leftmargin=1.4em}
\usepackage{fancyhdr}
\pagestyle{fancy}
\fancyhf{}
\fancyhead[L]{\small El problema gravitatorio de $N$ cuerpos}
\fancyhead[R]{\small J. García Izquierdo}
\fancyfoot[C]{\small\thepage}
\renewcommand{\headrulewidth}{0.3pt}
\setlength{\parskip}{0.35em}

\graphicspath{{figures/}}
\input{numbers.tex}

\newcommand{\vect}[1]{\mathbf{#1}}
\newcommand{\dd}{\mathrm{d}}
\newcommand{\repo}{\href{https://github.com/jesus-garcia-izquierdo/N-body-problem}{github.com/jesus-garcia-izquierdo/N-body-problem}}

\begin{document}

% ---------------------------------------------------------------------------
\begin{titlepage}
\centering
\vspace*{2.5cm}
{\Large\bfseries Jesús García Izquierdo - {\urlstyle{same}\hypersetup{urlcolor=black}\url{https://github.com/jesus-garcia-izquierdo/N-body-problem}}\par}
\vspace{0.6cm}
{\LARGE\bfseries El problema gravitatorio de $N$ cuerpos\par}
\vspace{0.25cm}
{\large Integración simpléctica, leyes de conservación, estabilidad y caos\par}
\vspace{1.5cm}
{\normalsize Octubre de 2026\par}
\vspace{1.2cm}
{\small Código Python, vídeos y datos: \repo\par}
\vspace{2cm}
\begin{minipage}{0.86\textwidth}
\small
\textbf{Resumen.} Estudiamos el movimiento de $N$ masas puntuales que solo interaccionan
mediante la gravedad newtoniana. Tras repasar la física (ecuaciones del movimiento,
estructura hamiltoniana, leyes de conservación, teorema del virial, estabilidad lineal y
exponentes de Lyapunov), integramos numéricamente las ecuaciones con el método de
Verlet-velocidad y otros cuatro esquemas, mediante una implementación en Python
totalmente vectorizada. Se analizan ocho experimentos: el sistema Sol--Tierra--Luna, un
sistema de cuatro cuerpos con un cuerpo central masivo, un sistema caótico de seis
cuerpos, una comparación de integradores, la coreografía en forma de ocho, la estabilidad
del triángulo equilátero de Lagrange, el Sistema Solar y el colapso gravitatorio de un
cúmulo estelar. La energía se conserva con un error de \SemErrE\ (Sol--Tierra--Luna, dos
años), los momentos lineal y angular hasta la precisión de la máquina, los órdenes de
convergencia medidos coinciden con la teoría, la tasa de inestabilidad del triángulo de
Lagrange concuerda con la teoría lineal en tres cifras (\LagSigma\ frente a
\LagSigmaTh), la tercera ley de Kepler se recupera con pendiente \SolarSlope\ y el
sistema de seis cuerpos tiene un tiempo de Lyapunov de solo \ChaosLyapDays\ días.
\end{minipage}
\vfill
{\footnotesize Este informe amplía y corrige el trabajo de 2023
\emph{Problema de los n-cuerpos} de G.~Haro Pérez y J.~García Izquierdo
(implementación en MATLAB).\par}
\end{titlepage}

\tableofcontents
\newpage

% ===========================================================================
\section{Introducción}
% ===========================================================================

El problema de $N$ cuerpos consiste en determinar el movimiento de $N$ masas que se
atraen gravitatoriamente. Para $N=2$ lo resolvió Newton: los cuerpos describen cónicas.
Para $N\ge3$ no existe una solución general en forma cerrada. Bruns (1887) demostró que no
hay más integrales algebraicas que las diez clásicas, y Poincaré (1890) mostró que el
movimiento puede ser extremadamente sensible a las condiciones iniciales, el fenómeno que
hoy llamamos \emph{caos}. Sundman (1912) encontró una serie de potencias convergente para
$N=3$, pero converge tan despacio que carece de utilidad práctica.

La única vía general es, por tanto, la integración numérica. Sin embargo, una solución
numérica solo es fiable si respeta lo que la dinámica exacta garantiza: la conservación de
la energía, del momento lineal y del momento angular, y relaciones estadísticas como el
teorema del virial. Este informe tiene tres objetivos:
\begin{enumerate}
  \item presentar la física necesaria para interpretar una simulación (sección~\ref{sec:model});
  \item explicar por qué un integrador \emph{simpléctico} como Verlet-velocidad es la
        herramienta adecuada y cuantificar su precisión (sección~\ref{sec:numerics});
  \item analizar ocho experimentos que van del Sistema Solar al caos y a la formación de
        cúmulos estelares (sección~\ref{sec:results}).
\end{enumerate}
Todos los valores numéricos del texto los produce el código del repositorio; el script
\texttt{tools/build\_documents.py} regenera todas las figuras y todos los números.

% ===========================================================================
\section{Modelo físico}\label{sec:model}
% ===========================================================================

\subsection{Ecuaciones del movimiento}

Sea $m_i$ la masa del cuerpo $i$, $\vect r_i(t)$ su posición y $\vect v_i(t)$ su velocidad.
La ley de gravitación de Newton junto con el principio de superposición da
\begin{equation}
  \ddot{\vect r}_i = G\sum_{j\ne i} m_j\,\frac{\vect r_j-\vect r_i}{|\vect r_j-\vect r_i|^{3}},
  \qquad i=1,\dots,N,
  \label{eq:eom}
\end{equation}
con $G=\SI{6.67430e-11}{m^3.kg^{-1}.s^{-2}}$. Son $3N$ ecuaciones diferenciales ordinarias
de segundo orden acopladas, o $6N$ de primer orden para $(\vect r_i,\vect v_i)$. Dadas las
posiciones y velocidades en $t=0$, la solución existe y es única mientras no choquen dos
cuerpos.

\subsection{Estructura hamiltoniana}

Escribiendo la energía en función de las posiciones $\vect q_i=\vect r_i$ y los momentos
$\vect p_i=m_i\vect v_i$,
\begin{equation}
  H(\vect q,\vect p)=\underbrace{\sum_i\frac{|\vect p_i|^2}{2m_i}}_{T(\vect p)}
  \;\underbrace{-\sum_{i<j}\frac{G m_i m_j}{|\vect q_i-\vect q_j|}}_{V(\vect q)},
  \label{eq:hamiltonian}
\end{equation}
las ecuaciones de Hamilton $\dot{\vect q}_i=\partial H/\partial\vect p_i$,
$\dot{\vect p}_i=-\partial H/\partial\vect q_i$ reproducen la ecuación~\eqref{eq:eom}. Dos
propiedades del flujo exacto guiarán la elección del integrador: conserva el volumen (de
hecho, la 2-forma simpléctica $\sum_i \dd\vect q_i\wedge \dd\vect p_i$) en el espacio de
fases de $6N$ dimensiones (teorema de Liouville) y es reversible en el tiempo. Además,
$H=T(\vect p)+V(\vect q)$ es \emph{separable}, lo que hace posible la descomposición de
Verlet (sección~\ref{sec:verlet}).

\subsection{Leyes de conservación}

Por el teorema de Noether, las simetrías de un sistema aislado implican cantidades
conservadas:
\begin{align}
  \text{traslaciones} &\;\Rightarrow\; \vect P=\sum_i m_i\vect v_i, \\
  \text{rotaciones}    &\;\Rightarrow\; \vect L=\sum_i \vect r_i\times m_i\vect v_i, \\
  \text{traslaciones temporales}  &\;\Rightarrow\; E=K+U,\quad K=\tfrac12\sum_i m_i v_i^2,\quad
                        U=-\sum_{i<j}\frac{G m_i m_j}{r_{ij}} .
  \label{eq:energy}
\end{align}
Las demostraciones indican qué debe respetar un método numérico. $\dot{\vect P}=0$ se
sigue de $\vect F_{ij}=-\vect F_{ji}$ (tercera ley de Newton), de modo que cualquier esquema
construido con fuerzas por pares exactamente antisimétricas conserva $\vect P$ salvo
errores de redondeo. $\dot{\vect L}=0$ se debe a que las fuerzas son centrales,
$\vect F_{ij}\parallel\vect r_i-\vect r_j$. La energía potencial cuenta cada par una sola
vez; sumar sobre todos los pares ordenados y dividir entre dos, como en el informe
original, es equivalente.

Como $\vect P$ es constante, el centro de masas (CM) se mueve uniformemente con
$\vect V_\text{cm}=\vect P/M$. Siempre que es posible trabajamos en el sistema del CM
($\vect P=0$): el sistema no se desplaza y la energía cinética pierde el término trivial
$\tfrac12 M V_\text{cm}^2$, imprescindible para el teorema del virial.

\subsection{El teorema del virial}\label{sec:virial}

Definimos el virial $\mathcal G=\sum_i\vect p_i\cdot\vect r_i$. Su derivada es
$\dot{\mathcal G}=2K+\sum_i\vect F_i\cdot\vect r_i$. Como $U$ es homogénea de grado $-1$ en
las posiciones, el teorema de Euler da $\sum_i\vect r_i\cdot\nabla_i U=-U$, es decir,
$\sum_i\vect F_i\cdot\vect r_i=U$. Promediando durante un tiempo $\tau$,
\begin{equation}
  \frac{\mathcal G(\tau)-\mathcal G(0)}{\tau}=2\langle K\rangle+\langle U\rangle
  \;\xrightarrow[\tau\to\infty]{}\;0
  \quad\Longrightarrow\quad
  Q\equiv\frac{2\langle K\rangle}{-\langle U\rangle}=1,
  \label{eq:virial}
\end{equation}
siempre que el movimiento esté acotado. Hacen falta, por tanto, tres condiciones: un
sistema ligado ($E<0$) sin cuerpos que escapen, $K$ medida en el sistema del CM y un tiempo
de promediado largo. En equilibrio, $E=\tfrac12\langle U\rangle=-\langle K\rangle$.

\subsection{El problema de dos cuerpos}

Con $\vect R$ el CM y $\vect r=\vect r_2-\vect r_1$, el problema de dos cuerpos se reduce
al de un cuerpo de masa reducida $\mu=m_1m_2/(m_1+m_2)$ en un campo central. La órbita es
una cónica; para órbitas ligadas usaremos repetidamente
\begin{equation}
  v^2=GM\left(\frac{2}{r}-\frac1a\right)\ \text{(vis-viva)},\qquad
  T^2=\frac{4\pi^2}{GM}\,a^3\ \text{(tercera ley de Kepler)},
\end{equation}
donde $a$ es el semieje mayor. Una órbita circular tiene $v_c=\sqrt{GM/r}$.

\subsection{Estabilidad, caos y exponentes de Lyapunov}\label{sec:chaos}

Sea $\boldsymbol\delta(t)$ la separación en el espacio de fases entre dos soluciones que
parten a una distancia $\delta_0$. Para casi cualquier dirección inicial,
\begin{equation}
  \lambda_\text{máx}=\lim_{t\to\infty}\lim_{\delta_0\to0}\frac1t\ln\frac{|\boldsymbol\delta(t)|}{\delta_0}
\end{equation}
existe y se denomina \emph{exponente máximo de Lyapunov}. El movimiento regular tiene
$\lambda_\text{máx}=0$ (las separaciones crecen, como mucho, de forma polinómica); el
movimiento caótico tiene $\lambda_\text{máx}>0$, y los errores se multiplican por $e$ cada
tiempo de Lyapunov $t_L=1/\lambda_\text{máx}$. Con una incertidumbre inicial $\delta_0$ y
una tolerancia $\Delta$, solo se puede predecir hasta
$t_\text{pred}\approx t_L\ln(\Delta/\delta_0)$: mejorar mil veces las medidas solo aporta
$\ln 1000\approx7$ tiempos de Lyapunov más.

Estimamos $\lambda_\text{máx}$ con el método de Benettin: junto a la trayectoria de
referencia se integra una trayectoria ``sombra''; cada intervalo $\tau$ se mide su
separación $d_k$, se acumula $\ln(d_k/\delta_0)$ y la sombra se reescala a distancia
$\delta_0$ en la misma dirección, de modo que $\lambda\approx t^{-1}\sum_k\ln(d_k/\delta_0)$.
Posiciones y velocidades se adimensionalizan con el radio cuadrático medio $L$ y el tiempo
$L/v_\text{rms}$.

La \emph{estabilidad lineal} es la versión local de la misma idea: una solución es
linealmente inestable si las perturbaciones pequeñas crecen exponencialmente. La solución
equilátera de Lagrange (sección~\ref{sec:lagrange}) es un ejemplo clásico en el que la tasa
de crecimiento puede calcularse analíticamente y compararse con la simulación.

\subsection{Suavizado y unidades}

En cúmulos estelares, donde las partículas representan muchas estrellas y los encuentros
cercanos entre dos cuerpos no interesan, sustituimos $1/r\to1/\sqrt{r^2+\varepsilon^2}$
(suavizado de Plummer). La fuerza sigue siendo el gradiente del potencial suavizado, de modo
que la energía suavizada también se conserva. Los problemas abstractos usan unidades de
$N$ cuerpos, $G=M=1$, en las que todas las magnitudes son de orden unidad; las escalas
físicas se recuperan teniendo en cuenta que el tiempo escala como $\sqrt{L^3/GM}$.

% ===========================================================================
\section{Método numérico}\label{sec:numerics}
% ===========================================================================

\subsection{Verlet-velocidad}\label{sec:verlet}

Escribiendo $\vect a_n=\vect a(\vect r_n)$ para las aceleraciones de la
ecuación~\eqref{eq:eom}, un paso de longitud $h$ es
\begin{align}
  \vect r_{n+1}&=\vect r_n+h\,\vect v_n+\tfrac{h^2}{2}\,\vect a_n, \label{eq:verlet1}\\
  \vect a_{n+1}&=\vect a(\vect r_{n+1}), \\
  \vect v_{n+1}&=\vect v_n+\tfrac{h}{2}\left(\vect a_n+\vect a_{n+1}\right). \label{eq:verlet3}
\end{align}
La actualización de la posición es un desarrollo de Taylor de segundo orden y la de la
velocidad es la regla del trapecio para $\int\vect a\,\dd t$. Introduciendo
$\vect v_{n+1/2}=\vect v_n+\tfrac h2\vect a_n$, el mismo algoritmo se lee ``medio impulso --
deriva -- medio impulso''. Se trata de la descomposición de operadores
$e^{\frac h2 V}e^{hT}e^{\frac h2 V}$ del hamiltoniano~\eqref{eq:hamiltonian}: la deriva es
el flujo exacto de $T$ y los impulsos son el flujo exacto de $V$. Por tanto:
\begin{itemize}
  \item el método es \textbf{simpléctico} (composición de flujos hamiltonianos) y
        \textbf{reversible} (la composición es simétrica);
  \item por la fórmula de Baker--Campbell--Hausdorff, el error local es $\mathcal O(h^3)$ y
        el global $\mathcal O(h^2)$;
  \item conserva $\vect P$ y $\vect L$ exactamente (cada subpaso lo hace);
  \item necesita una evaluación de fuerzas por paso, porque $\vect a_{n+1}$ se reutiliza
        como $\vect a_n$ del paso siguiente (el código MATLAB original la calculaba dos veces).
\end{itemize}

\textbf{Por qué importa que sea simpléctico.} El análisis de error regresivo muestra que un
método simpléctico de orden $p$ sigue \emph{exactamente} la dinámica de un hamiltoniano
modificado o ``sombra'' $\tilde H=H+h^pH_p+\dots$, salvo términos exponencialmente
pequeños. La trayectoria numérica conserva $\tilde H$, así que el error en $H$ permanece
acotado por $\mathcal O(h^p)$ durante tiempos exponencialmente largos en lugar de derivar.
Es la propiedad más importante para integraciones orbitales largas.

\subsection{Otros integradores}

Para mostrar el efecto del orden y de la estructura se han implementado otros cuatro
esquemas (tabla~\ref{tab:integrators}). Euler explícito multiplica las áreas del espacio de
fases por $1+\omega^2h^2$ en cada paso para un oscilador armónico, de modo que la energía
crece exponencialmente; intercambiar dos líneas (Euler simpléctico) hace que ese factor sea
exactamente uno. El método de Yoshida compone tres pasos de Verlet de longitudes
$w_1h,\,w_0h,\,w_1h$ con $w_1=1/(2-2^{1/3})$ y $w_0=-2^{1/3}/(2-2^{1/3})$ (obsérvese el paso
hacia atrás en el tiempo), lo que cancela el término de error $h^3$. RK4 es de cuarto orden
pero no es simpléctico.

\begin{table}[H]
\centering
\caption{Integradores implementados en \texttt{nbody/integrators.py}.}
\label{tab:integrators}
\begin{tabular}{lcccc}
\toprule
Método & Orden & Simpléctico & Reversible & Fuerzas por paso\\
\midrule
Euler explícito    & 1 & no & no & 1\\
Euler simpléctico  & 1 & sí & no & 1\\
Verlet-velocidad   & 2 & sí & sí & 1\\
Yoshida-4          & 4 & sí & sí & 3\\
Runge--Kutta 4     & 4 & no & no & 4\\
\bottomrule
\end{tabular}
\end{table}

\subsection{Implementación}

El código es un pequeño paquete de Python (\texttt{nbody}) basado en NumPy. El estado se
guarda en arrays de forma $(N,3)$ (cuerpo $\times$ componente) y una trayectoria en arrays
de forma $(T,N,3)$ (instante $\times$ cuerpo $\times$ componente). Las aceleraciones se
calculan sin ningún bucle de Python sobre los cuerpos: el \emph{broadcasting} construye el
tensor $(N,N,3)$ de separaciones $\vect r_j-\vect r_i$, la diagonal se excluye haciendo
$r_{ii}^2=\infty$ y una suma de Einstein contrae el índice $j$. El coste es
$\mathcal O(N^2)$ por paso, como en cualquier código de suma directa. La memoria de la
trayectoria se reserva una sola vez (la versión MATLAB hacía crecer los arrays en cada paso,
un coste $\mathcal O(T^2)$). En un portátil, el sistema Sol--Tierra--Luna
($6{,}3\times10^4$ pasos) tarda alrededor de un segundo y el Sistema Solar
($1{,}2\times10^5$ pasos), unos pocos segundos.

\subsection{Diagnósticos y elección del paso}

La conservación se mide \emph{respecto al valor inicial}: $|E(t)-E_0|/|E_0|$ para la
energía y, para los vectores $\vect P$ y $\vect L$, la norma del cambio dividida por las
escalas naturales $\max_t\sum_i m_i|\vect v_i|$ y $\max_t\sum_i m_i|\vect r_i||\vect v_i|$
(necesarias porque $\vect P=0$ en el sistema del CM). El paso temporal lo fija el tiempo
dinámico más corto del sistema, con la regla práctica $h\lesssim t_\text{din}/100$, y todos
los resultados se han comprobado reduciendo $h$ a la mitad.

% ===========================================================================
\section{Resultados}\label{sec:results}
% ===========================================================================

\subsection{El sistema Sol--Tierra--Luna}

Los cuerpos parten en el plano de la eclíptica, con la Tierra en el afelio y la Luna en el
apogeo (tabla~\ref{tab:sem}). El periodo más corto es el de la Luna, 27,3 días, de modo que
$h=\SI{1000}{s}$ proporciona unos 2400 pasos por órbita lunar.

\begin{table}[H]
\centering
\caption{Condiciones iniciales del sistema Sol--Tierra--Luna (antes de pasar al sistema del
CM). Tiempo de integración: 2 años; $h=\SI{1000}{s}$.}
\label{tab:sem}
\begin{tabular}{lccc}
\toprule
 & Sol & Tierra & Luna\\
\midrule
Masa (kg) & \num{1.98847e30} & \num{5.9722e24} & \num{7.342e22}\\
$x_0$ (m) & 0 & \num{1.52098232e11} (afelio) & $x_\oplus+\num{4.055e8}$ (apogeo)\\
$v_{y,0}$ (m/s) & 0 & \num{29290} & $v_\oplus+\num{964}$\\
\bottomrule
\end{tabular}
\end{table}

La figura~\ref{fig:sem-orbits} muestra las órbitas. A escala heliocéntrica la órbita de la
Luna es invisible (la distancia Tierra--Luna es 1/375 de la distancia Sol--Tierra), por lo
que se representa también respecto a la Tierra: la elipse lunar \emph{precesa} de forma
visible, un efecto directo de la perturbación de marea del Sol que un modelo de dos cuerpos
no recogería.

\begin{figure}[H]
\centering
\includegraphics[width=\textwidth]{sem_orbits}
\caption{Sistema Sol--Tierra--Luna durante dos años. Izquierda: órbitas alrededor del centro
de masas (la Tierra y la Luna se superponen a esta escala). Derecha: la Luna respecto a la
Tierra; su órbita precesa por la perturbación del Sol.}
\label{fig:sem-orbits}
\end{figure}

Las leyes de conservación se muestran en la figura~\ref{fig:sem-cons}. El error de energía
oscila con el periodo orbital y se mantiene por debajo de \SemErrE: no hay deriva secular,
como cabe esperar de un método simpléctico. Los momentos lineal y angular se conservan con
errores de \SemErrP\ y \SemErrL, respectivamente, es decir, hasta la precisión de la
máquina; el lento crecimiento es la acumulación de errores de redondeo. Los promedios
temporales dan $\langle K\rangle=\SemKmean\,\mathrm{J}$ y
$-\tfrac12\langle U\rangle=\SemUhalf\,\mathrm{J}$, de modo que $Q=\SemVirial$: el teorema
del virial se cumple con un error inferior a $10^{-4}$.

\begin{figure}[H]
\centering
\includegraphics[width=\textwidth]{sem_conservation}
\caption{Leyes de conservación en la simulación Sol--Tierra--Luna: (a) energías cinética,
potencial y total; (b) error relativo de la energía, acotado y oscilante; (c) errores
relativos de $\vect P$ y $\vect L$, al nivel del redondeo.}
\label{fig:sem-cons}
\end{figure}

Si el Sol parte en reposo (como en el informe original), el momento total no es nulo y el
CM se mueve a $V_\text{cm}=\SemVcm$~m/s. Aquí es despreciable porque el Sol domina la
masa, pero en general sesga $\langle K\rangle$ y, por tanto, $Q$.

\subsection{Cuatro cuerpos alrededor de un cuerpo central masivo}

Este sistema reproduce la sección~4.2 del informe original (tabla~\ref{tab:four}): un
cuerpo central de \SI{2e40}{kg} (unas $10^{10}$ masas solares, comparable a un núcleo
galáctico), dos estrellas de aproximadamente una masa solar en el plano $xy$ y un cuerpo
ligero en una órbita polar.

\begin{table}[H]
\centering
\caption{Condiciones iniciales del sistema de cuatro cuerpos ($d=\SI{1.496e14}{m}$,
$M_\odot=\SI{1.9891e30}{kg}$). Tiempo de integración: 15 años; $h=\SI{1000}{s}$.}
\label{tab:four}
\begin{tabular}{lcccc}
\toprule
 & Cuerpo 1 & Cuerpo 2 & Cuerpo 3 & Cuerpo 4\\
\midrule
Masa (kg) & $2{,}02\,M_\odot$ & $0{,}98\,M_\odot$ & \num{2e40} & \num{2e20}\\
$\vect r_0$ (m) & $(0,0,0)$ & $(d,0,0)$ & $(d/2,0,0)$ & $(d/2,0,10^{14})$\\
$\vect v_0$ & $(0,-\sqrt{Gm_3/d},0)$ & $(0,\sqrt{Gm_3/d},0)$ & $\vect 0$ & $(-\sqrt{Gm_3/10^{14}},0,0)$\\
\bottomrule
\end{tabular}
\end{table}

El informe original pretendía órbitas circulares, pero las velocidades usan la distancia
$d$, que es el \emph{doble} de la distancia real $r=d/2$ al cuerpo central. Entonces
$v^2=GM/(2r)$ y la ecuación vis-viva da $a=2r/3$: las órbitas son elipses de excentricidad
$e=0{,}5$ cuyo apocentro es el punto de partida. La simulación mide $e=\FourEcc$, lo que
confirma el análisis. El cuerpo 4 parte a la distancia correcta y describe una circunferencia.

\begin{figure}[H]
\centering
\includegraphics[width=0.82\textwidth]{four_orbits}
\caption{Órbitas del sistema de cuatro cuerpos durante 15 años (unas 250 órbitas de los
cuerpos 1 y 2). Las elipses son cerradas y estables; el cuerpo masivo apenas se mueve.}
\label{fig:four-orbits}
\end{figure}

\begin{figure}[H]
\centering
\includegraphics[width=\textwidth]{four_conservation}
\caption{Leyes de conservación en el sistema de cuatro cuerpos. El error de energía es
máximo en los pasos por el pericentro, donde el paso temporal es relativamente mayor, pero no
crece.}
\label{fig:four-cons}
\end{figure}

Con órbitas excéntricas, el error de energía alcanza su máximo en cada paso por el
pericentro, donde la escala temporal orbital es más corta, y llega a \FourErrE\
(figura~\ref{fig:four-cons}); aun así, su envolvente es plana durante 15 años. $\vect P$ y
$\vect L$ se conservan con errores de \FourErrP\ y \FourErrL, y el cociente del virial es
$Q=\FourVirial$.

\subsection{Caos en un sistema de seis cuerpos}

Seis cuerpos de \SI{1e10}{kg} parten de los semiejes $\pm x$, $\pm y$, $\pm z$ a
\SI{1e5}{m}, cada uno con la velocidad ($\approx\SI{2.1}{mm/s}$) de una órbita circular
alrededor del siguiente (tabla~\ref{tab:six}). Integramos cuatro versiones durante 10 años:
la de referencia y tres en las que una única magnitud cambia un 0,1\,\%.

\begin{table}[H]
\centering
\caption{Sistema de seis cuerpos ($u=\SI{1.5e-3}{m/s}$). Tiempo de integración: 10 años;
$h=\SI{1000}{s}$.}
\label{tab:six}
\begin{tabular}{lcccccc}
\toprule
 & 1 & 2 & 3 & 4 & 5 & 6\\
\midrule
$\vect r_0/\SI{e5}{m}$ & $(1,0,0)$ & $(0,1,0)$ & $(0,0,1)$ & $(-1,0,0)$ & $(0,-1,0)$ & $(0,0,-1)$\\
$\vect v_0/u$ & $(1,1,0)$ & $(0,1,1)$ & $(-1,0,1)$ & $(-1,-1,0)$ & $(0,-1,-1)$ & $(1,0,-1)$\\
\bottomrule
\end{tabular}
\end{table}

\begin{figure}[H]
\centering
\includegraphics[width=\textwidth]{chaos_orbits}
\caption{Sistema de seis cuerpos: simulación de referencia y tres simulaciones con un cambio
del 0,1\,\% en $x_1$, $v_{y,2}$ y $m_3$. Los cuerpos se emparejan en binarias y algunos son
expulsados, de forma distinta en cada caso.}
\label{fig:chaos-orbits}
\end{figure}

La figura~\ref{fig:chaos-orbits} muestra que las configuraciones finales son
cualitativamente distintas. La figura~\ref{fig:chaos-div} lo cuantifica: la separación
respecto a la referencia crece exponencialmente (una recta en escala logarítmica) tras el
primer encuentro cercano, en torno a $t\approx2$ años, hasta saturar al tamaño del sistema.

La curva más instructiva es la última: la \emph{misma} condición inicial integrada con $h$
y con $h/2$ también diverge exponencialmente, y alcanza \ChaosHalfFinal~m al cabo de 10
años. El error de truncamiento actúa como una perturbación más, de modo que, pasados unos
pocos tiempos de Lyapunov, ninguna de las dos simulaciones reproduce la trayectoria exacta.
Esto no invalida la simulación: $E$, $\vect P$ y $\vect L$ se siguen conservando
(\ChaosErrE, \ChaosErrP, \ChaosErrL), y el lema de \emph{shadowing} garantiza que la órbita
numérica permanece cerca de \emph{alguna} órbita exacta con condiciones iniciales
ligeramente distintas. Las propiedades estadísticas son fiables; las trayectorias
individuales, no.

\begin{figure}[H]
\centering
\includegraphics[width=\textwidth]{chaos_divergence}
\caption{Separación $\delta r(t)$ entre cada simulación perturbada y la de referencia. La
curva amarilla compara la misma condición inicial integrada con $h$ y con $h/2$.}
\label{fig:chaos-div}
\end{figure}

El método de Benettin (figura~\ref{fig:chaos-lyap}) da
$\lambda_\text{máx}\approx\ChaosLambdaYr~\text{año}^{-1}$ al cabo de 10 años, es decir, un
tiempo de Lyapunov de unos \ChaosLyapDays\ días. La estimación salta tras el primer
encuentro cercano y todavía deriva lentamente, señal de que la fase caótica comienza en
$t\approx2$~años. En diez años, un error inicial se amplifica en unos \ChaosAmpExp\ órdenes
de magnitud: ni siquiera unas condiciones iniciales con precisión doble (error relativo
$10^{-16}$) permitirían una predicción determinista.

\begin{figure}[H]
\centering
\includegraphics[width=\textwidth]{chaos_lyapunov}
\caption{Estimación acumulada del exponente máximo de Lyapunov (método de Benettin,
renormalización cada 20 pasos, $\delta_0=10^{-8}$). La línea discontinua es el promedio del
último 20\,\% de la simulación.}
\label{fig:chaos-lyap}
\end{figure}

\subsection{Comparación de integradores}

Para aislar el efecto del integrador usamos la órbita en forma de ocho de la
sección~\ref{sec:eight}, con $G=m=1$ y periodo $T\approx6{,}326$. La
figura~\ref{fig:int-energy} muestra el error máximo de energía por periodo durante 200
periodos con un paso deliberadamente grueso, $h=T/300$. Euler explícito gana energía en
cada paso y el sistema se desintegra en menos de diez periodos. RK4 empieza con un error
diminuto, pero deriva \emph{linealmente}. Los tres esquemas simplécticos no muestran ninguna
deriva, como predice el argumento del hamiltoniano sombra; a igualdad de coste, Verlet es
mucho más robusto de lo que sugiere su orden formal, y Yoshida-4 combina orden alto con
error acotado.

\begin{figure}[H]
\centering
\includegraphics[width=\textwidth]{integrators_energy}
\caption{Error de energía a largo plazo (máximo por periodo) con $h=T/300$ en la órbita en
forma de ocho. Solo los métodos no simplécticos muestran una deriva secular.}
\label{fig:int-energy}
\end{figure}

La figura~\ref{fig:int-conv} representa el error en la posición tras un periodo frente al
número de evaluaciones de fuerza, que es el coste real. Las pendientes ajustadas dan órdenes
\OrderEuler, \OrderSympEuler, \OrderVerlet, \OrderYoshida\ y \OrderRK, que coinciden con
los valores teóricos 1, 1, 2, 4 y 4.

\begin{figure}[H]
\centering
\includegraphics[width=\textwidth]{integrators_convergence}
\caption{Orden de convergencia: error en la posición tras un periodo frente al número de
evaluaciones de fuerza. Las pendientes se ajustan con los tres pasos más pequeños.}
\label{fig:int-conv}
\end{figure}

\subsection{La coreografía en forma de ocho}\label{sec:eight}

Moore (1993) encontró numéricamente, y Chenciner y Montgomery (2000) demostraron, que tres
masas iguales pueden perseguirse a lo largo de una única curva en forma de ocho, separadas
un tercio de periodo. Las condiciones iniciales son
$\vect r_1=-\vect r_2=(-0{,}97000436;\;0{,}24308753)$, $\vect r_3=\vect 0$,
$\vect v_3=(-0{,}93240737;\;-0{,}86473146)$, $\vect v_1=\vect v_2=-\vect v_3/2$. El momento
angular total es exactamente nulo aunque los cuerpos giran: los dos lóbulos se recorren en
sentidos opuestos.

Con $h=T/6000$ (figura~\ref{fig:eight}), los cuerpos vuelven a sus posiciones iniciales con
un error de \EightRetOne\ tras un periodo y de \EightRetTen\ tras diez (los propios datos
iniciales se dan con ocho cifras). El error de energía está acotado en \EightErrE\ y
$Q=\EightVirial$.

\begin{figure}[H]
\centering
\includegraphics[width=\textwidth]{eight_orbit}
\caption{Coreografía en forma de ocho durante diez periodos: las tres órbitas coinciden.
Derecha: error relativo de la energía, acotado.}
\label{fig:eight}
\end{figure}

Cambiar una componente de la velocidad un 0,1\,\% produce una separación que crece solo
lentamente y \emph{sin} la tendencia exponencial del sistema de seis cuerpos
(figura~\ref{fig:eight-stab}), y que llega a \EightDivFinal\ tras diez periodos. La órbita
está rodeada de toros KAM invariantes: es estable.

\begin{figure}[H]
\centering
\includegraphics[width=\textwidth]{eight_stability}
\caption{Separación entre la órbita en forma de ocho y una copia con $v_{x,1}$ modificada un
0,1\,\%: crecimiento lento y no exponencial.}
\label{fig:eight-stab}
\end{figure}

\subsection{Estabilidad del triángulo de Lagrange}\label{sec:lagrange}

Lagrange (1772) demostró que tres cuerpos en los vértices de un triángulo equilátero de lado
$d$, girando rígidamente con $\omega^2=GM/d^3$, forman una solución exacta para cualesquiera
masas. Linealizando en el sistema de referencia que rota, los exponentes de las
perturbaciones cumplen
\begin{equation}
  \lambda^4+\omega^2\lambda^2+\tfrac{27}{4}\,\beta\,\omega^4=0,\qquad
  \beta=\frac{m_1m_2+m_2m_3+m_3m_1}{M^2}.
\end{equation}
Las raíces son imaginarias puras (oscilaciones estables) solo si $27\beta<1$, que es el
criterio de Routh. Para el Sol, Júpiter y un asteroide troyano ($\beta\approx10^{-3}$) la
configuración es estable, y por eso los asteroides troyanos se agrupan en los puntos L4 y
L5 de Júpiter. Para masas iguales, $\beta=1/3$, las raíces cumplen
$\lambda^2=\omega^2(-1\pm i\sqrt8)/2$ y la tasa de crecimiento es
$\operatorname{Re}\lambda=\omega/\sqrt2$.

En la simulación no se añade ninguna perturbación a mano: basta el error de truncamiento de
Verlet (figura~\ref{fig:lagrange}). Con masas iguales, la deformación del triángulo crece
exponencialmente desde $10^{-6}$ hasta orden unidad, y un ajuste lineal de $\ln\delta$ da
$\sigma=\LagSigma$, frente al valor teórico $\LagSigmaTh$. Con las masas Sol--Júpiter--troyano,
la deformación se mantiene por debajo de \LagTrojanMax. La pequeña deformación periódica del
principio no es una inestabilidad: Verlet resuelve exactamente un hamiltoniano sombra cuya
órbita circular tiene un radio ligeramente distinto.

\begin{figure}[H]
\centering
\includegraphics[width=\textwidth]{lagrange}
\caption{Deformación del triángulo de Lagrange (mayor cambio relativo de un lado). Masas
iguales: crecimiento exponencial con la tasa que predice la teoría lineal.
Sol--Júpiter--troyano: acotada.}
\label{fig:lagrange}
\end{figure}

\subsection{El Sistema Solar y la tercera ley de Kepler}

El Sol y los ocho planetas se inicializan con sus masas, semiejes mayores, excentricidades e
inclinaciones reales (tabla~\ref{tab:planets}). Cada planeta parte del perihelio con
velocidad $v_p=\sqrt{GM_\odot(1+e)/(a(1-e))}$ y una fase orbital aleatoria; su plano orbital
se inclina girando alrededor de la recta Sol--perihelio. Son condiciones iniciales
ilustrativas, no las efemérides de una fecha concreta. El paso $h=0{,}5$~días lo fija
Mercurio (176 pasos por órbita), y la simulación abarca una órbita de Neptuno, 165 años.

\begin{table}[H]
\centering
\caption{Parámetros planetarios y periodos orbitales medidos en la simulación (contando
vueltas del ángulo polar alrededor del Sol).}
\label{tab:planets}
\begin{tabular}{lccccc}
\toprule
Planeta & Masa (kg) & $a$ (UA) & $e$ & $i$ (°) & $T$ medido (años)\\
\midrule
Mercurio & \num{3.3011e23} & 0,387 & 0,206 & 7,00 & \PeriodMercury\\
Venus    & \num{4.8675e24} & 0,723 & 0,007 & 3,39 & \PeriodVenus\\
Tierra   & \num{5.9722e24} & 1,000 & 0,017 & 0,00 & \PeriodEarth\\
Marte    & \num{6.4171e23} & 1,524 & 0,093 & 1,85 & \PeriodMars\\
Júpiter  & \num{1.8982e27} & 5,203 & 0,049 & 1,30 & \PeriodJupiter\\
Saturno  & \num{5.6834e26} & 9,555 & 0,057 & 2,49 & \PeriodSaturn\\
Urano    & \num{8.6810e25} & 19,22 & 0,046 & 0,77 & \PeriodUranus\\
Neptuno  & \num{1.02413e26} & 30,11 & 0,011 & 1,77 & \PeriodNeptune\\
\bottomrule
\end{tabular}
\end{table}

\begin{figure}[H]
\centering
\includegraphics[width=\textwidth]{solar_orbits}
\caption{Sistema Solar: planetas interiores durante dos años (izquierda) y el sistema
completo durante una órbita de Neptuno (derecha).}
\label{fig:solar}
\end{figure}

Un ajuste log-log de los periodos medidos frente a los semiejes mayores
(figura~\ref{fig:kepler}) da una pendiente de \SolarSlope, frente al valor $3/2$ de la
tercera ley de Kepler. Las pequeñas desviaciones de los planetas exteriores son reales:
Júpiter y Saturno se perturban mutuamente y perturban a sus vecinos, y para Urano y Neptuno
una sola órbita excéntrica es demasiado corta para promediar la velocidad angular. En 165
años, la energía se conserva con un error de \SolarErrE, los momentos con \SolarErrP\ y
\SolarErrL, y $Q=\SolarVirial$.

\begin{figure}[H]
\centering
\includegraphics[width=0.75\textwidth]{solar_kepler}
\caption{Tercera ley de Kepler: periodos orbitales medidos en la simulación frente al
semieje mayor, junto con la recta $T=a^{3/2}$ (años y UA).}
\label{fig:kepler}
\end{figure}

\subsection{Cúmulos estelares y el teorema del virial}

Por último, seguimos 400 ``estrellas'' de igual masa en unidades de $N$ cuerpos, con
suavizado $\varepsilon=0{,}05$ y $h=0{,}002$. Una \emph{esfera de Plummer}, muestreada a
partir de su función de distribución de equilibrio, nace virializada y se mantiene cerca de
$Q=1$ (promedio a tiempos largos \PlummerQlate), con fluctuaciones debidas a que $N$ es
finito.

Un \emph{colapso frío} parte de una esfera homogénea con todas las estrellas en reposo
($Q=0$). La gravedad la hace colapsar en aproximadamente un tiempo de caída libre
$t_{ff}=\frac\pi2\sqrt{R^3/2GM}\approx1{,}11$; en la compresión máxima, $Q$ sobrepasa la
unidad hasta \ColdQmax, el cúmulo rebota y en pocos $t_{ff}$ se asienta en un estado
cuasiestacionario con $\langle Q\rangle=\ColdQlate$ (figura~\ref{fig:cluster}). Esta
\emph{relajación violenta} (Lynden-Bell, 1967) ocurre en la escala de tiempo dinámica, sin
colisiones entre estrellas, y se cree que da forma a las galaxias elípticas. Los radios
lagrangianos muestran la formación de un núcleo denso mientras el 10\,\% exterior de la masa
se expande formando un halo, parte del cual escapa (figura~\ref{fig:collapse}). La energía
suavizada se conserva con un error de \ColdErrE\ pese a la violenta compresión.

\begin{figure}[H]
\centering
\includegraphics[width=\textwidth]{cluster_virial}
\caption{Izquierda: cociente del virial $Q=2K/|U|$ en el colapso frío y en la esfera de
Plummer. Derecha: radios lagrangianos que contienen el 10, el 50 y el 90\,\% de la masa
durante el colapso frío.}
\label{fig:cluster}
\end{figure}

\begin{figure}[H]
\centering
\includegraphics[width=\textwidth]{cluster_snapshots}
\caption{Instantáneas del colapso frío (proyección $xy$): esfera inicial, contracción,
compresión máxima y cúmulo final relajado con estrellas que escapan.}
\label{fig:collapse}
\end{figure}

\subsection{Vídeos}

El repositorio contiene además vídeos 3D (MP4) de estos experimentos, con cámara giratoria
y estelas que se desvanecen (\texttt{videos/}). La figura~\ref{fig:stills} muestra cuatro
fotogramas.

\IfFileExists{figures/video_stills.png}{
\begin{figure}[H]
\centering
\includegraphics[width=\textwidth]{video_stills.png}
\caption{Fotogramas de los vídeos: coreografía en forma de ocho, caos en el sistema de seis
cuerpos, el Sistema Solar durante una órbita de Neptuno y el colapso frío de un cúmulo
estelar.}
\label{fig:stills}
\end{figure}
}{}

% ===========================================================================
\section{Discusión}
% ===========================================================================

La tabla~\ref{tab:summary} recoge los errores de conservación de todas las simulaciones. Los
momentos lineal y angular se conservan hasta la precisión de la máquina en todos los casos,
como se espera de un método basado en fuerzas centrales por pares antisimétricas. El error
de energía está siempre acotado y su tamaño lo controla el cociente entre el paso temporal y
el tiempo dinámico más corto: es mayor en las órbitas excéntricas del sistema de cuatro
cuerpos y en el núcleo denso del cúmulo en colapso.

\begin{table}[H]
\centering
\caption{Errores máximos de conservación y cociente del virial en cada simulación
(Verlet-velocidad).}
\label{tab:summary}
\begin{tabular}{lcccc}
\toprule
Simulación & $|\Delta E/E_0|$ & $|\Delta\vect P|$ (rel.) & $|\Delta\vect L|$ (rel.) & $Q$\\
\midrule
Sol--Tierra--Luna (2 años) & \SemErrE & \SemErrP & \SemErrL & \SemVirial\\
Cuatro cuerpos (15 años)   & \FourErrE & \FourErrP & \FourErrL & \FourVirial\\
Seis cuerpos (10 años)     & \ChaosErrE & \ChaosErrP & \ChaosErrL & ---\\
Ocho (10 $T$)              & \EightErrE & --- & --- & \EightVirial\\
Sistema Solar (165 años)   & \SolarErrE & \SolarErrP & \SolarErrL & \SolarVirial\\
\bottomrule
\end{tabular}
\end{table}

El teorema del virial se cumple con un error de $10^{-4}$ o inferior en todos los sistemas
ligados y acotados. No se aplica al sistema de seis cuerpos, donde hay cuerpos que escapan,
ni a los primeros instantes del colapso frío, que es precisamente lo que convierte a este
último en una buena ilustración de la relajación hacia el equilibrio virial.

Respecto al trabajo original de 2023 conviene señalar tres correcciones. Primero, la masa de
la Luna de la tabla~1 de aquel informe (\SI{2e40}{kg}) era una errata por
\SI{7.342e22}{kg}. Segundo, la órbita terrestre tiene una excentricidad cercana a cero, no a
uno. Tercero, las órbitas del sistema de cuatro cuerpos no son circulares sino que tienen
$e=0{,}5$, como explica la ecuación vis-viva. Además, medir los errores respecto al valor
inicial (en lugar de restar la media) revela la lenta deriva de $\vect L$ por redondeo que
la media oculta.

Las principales limitaciones son las de cualquier código de paso fijo y suma directa. Los
encuentros cercanos (por ejemplo, el problema pitagórico de tres cuerpos de Burrau) exigen
pasos adaptativos, que deben diseñarse con cuidado para preservar la simplecticidad, o una
regularización. El coste $\mathcal O(N^2)$ limita el método a unos pocos miles de cuerpos;
más allá hacen falta códigos de árbol ($\mathcal O(N\log N)$) o métodos multipolares
rápidos.

% ===========================================================================
\section{Conclusiones}
% ===========================================================================

\begin{itemize}
  \item El método de Verlet-velocidad conserva los momentos lineal y angular hasta la
        precisión de la máquina y mantiene acotado el error de energía durante miles de
        órbitas, porque es simpléctico. Los métodos no simplécticos derivan, aunque sean de
        orden más alto.
  \item Los órdenes de convergencia medidos (\OrderVerlet\ para Verlet, \OrderYoshida\ para
        Yoshida-4) concuerdan con la teoría, y Yoshida-4 ofrece la mejor precisión por coste
        sin dejar de ser simpléctico.
  \item Se reproducen cuantitativamente resultados analíticos: la tercera ley de Kepler
        (pendiente \SolarSlope), la excentricidad que predice la ecuación vis-viva
        (\FourEcc), la tasa de inestabilidad lineal del triángulo de Lagrange (\LagSigma\
        frente a \LagSigmaTh) y el teorema del virial.
  \item El problema de $N$ cuerpos es, en general, caótico: el sistema de seis cuerpos tiene
        un tiempo de Lyapunov de unos \ChaosLyapDays\ días, e incluso el paso de integración
        cambia la trayectoria tras unos pocos tiempos de Lyapunov. Las predicciones a largo
        plazo solo tienen sentido estadístico, garantizado por las leyes de conservación y
        la propiedad de \emph{shadowing}.
\end{itemize}

% ===========================================================================
\appendix
\section{Cómo reproducir los resultados}
% ===========================================================================

\begin{verbatim}
git clone https://github.com/jesus-garcia-izquierdo/N-body-problem.git
cd N-body-problem
pip install -e .
python examples/01_sun_earth_moon.py      # ... hasta 08_cluster_virial.py
python tools/build_documents.py           # regenera este informe
\end{verbatim}

% ===========================================================================
\begin{thebibliography}{99}
\small
\bibitem{aarseth} S.~J. Aarseth, \emph{Gravitational N-Body Simulations}, Cambridge University Press (2003).
\bibitem{benettin} G.~Benettin, L.~Galgani, A.~Giorgilli, J.-M. Strelcyn, \emph{Meccanica} \textbf{15}, 9--30 (1980).
\bibitem{binney} J.~Binney, S.~Tremaine, \emph{Galactic Dynamics}, 2.ª ed., Princeton University Press (2008).
\bibitem{chenciner} A.~Chenciner, R.~Montgomery, \emph{Annals of Mathematics} \textbf{152}, 881--901 (2000).
\bibitem{goldstein} H.~Goldstein, C.~Poole, J.~Safko, \emph{Classical Mechanics}, 3.ª ed., Addison-Wesley (2002).
\bibitem{hairer} E.~Hairer, C.~Lubich, G.~Wanner, \emph{Geometric Numerical Integration}, 2.ª ed., Springer (2006).
\bibitem{haro} G.~Haro Pérez, J.~García Izquierdo, \emph{Problema de los n-cuerpos}, informe de asignatura (2023).
\bibitem{laskar} J.~Laskar, \emph{Nature} \textbf{338}, 237--238 (1989).
\bibitem{lyndenbell} D.~Lynden-Bell, \emph{Monthly Notices of the RAS} \textbf{136}, 101 (1967).
\bibitem{moore} C.~Moore, \emph{Physical Review Letters} \textbf{70}, 3675 (1993).
\bibitem{szebehely} V.~Szebehely, C.~F. Peters, \emph{Astronomical Journal} \textbf{72}, 876 (1967).
\bibitem{verlet} L.~Verlet, \emph{Physical Review} \textbf{159}, 98 (1967).
\bibitem{yoshida} H.~Yoshida, \emph{Physics Letters A} \textbf{150}, 262--268 (1990).
\end{thebibliography}

\end{document}
